\section{Introduction}

This sprint is dedicated to attack graph generation. The goal is to transform project artifacts into diagrams that show how data, endpoints and vulnerabilities are connected, so these diagrams highlight potential attack paths in the platform.

\section{Sprint Purpose}

The purpose of this sprint is to implement the Dashboard's attack-path and scenario generation features. The generator is an AI-assisted service that produces ephemeral attack scenarios (ordered steps) and a graph object (nodes and edges) for visualization in the Dashboard. Unlike static report diagrams, these graphs are runtime artifacts produced on-demand and returned to the frontend for interactive rendering.

\section{Sprint Backlog}

The backlog for this sprint focuses on the AI-driven scenario and risk-generation flows and the Dashboard UI integration:

\begin{table}[H]
  \centering
  \renewcommand{\arraystretch}{1.2}
  \begin{tabularx}{\textwidth}{|X|}
    \hline
    As an analyst, I can request an attack scenario for a specific vulnerability and receive structured steps and a graph object. \\\hline
    As an analyst, I can request a risk assessment for a completed scan and receive a summarized JSON assessment. \\\hline
    As the system, I must validate and normalize AI responses into a safe, consistent graph payload (nodes/edges). \\\hline
    As a frontend, the Dashboard should render the returned graph (React Flow) and allow exporting the scenario as JSON. \\\hline
    As an operator, the service should handle AI timeouts, malformed responses, and return useful error messages. \\\hline
  \end{tabularx}
  \caption{Sprint Backlog}
\end{table}

\section{Requirements Analysis}


\subsection{Generation Pipeline}

The generator implements a compact, AI-assisted runtime pipeline: the Dashboard requests a scenario for a selected vulnerability; the backend loads the vulnerability context, builds a controlled prompt, calls the configured AI gateway, extracts and validates the returned JSON, maps normalized steps into an ephemeral graph (nodes and edges), and returns the graph and ordered steps to the frontend for interactive rendering.

\begin{figure}[H]
  \centering
  \fbox{\parbox{0.9\textwidth}{\centering Placeholder: runtime pipeline (Vulnerability/Scan → AI gateway → parser → ephemeral graph)}}
  \caption{Attack graph runtime pipeline (placeholder)}
  \label{fig:graph-pipeline}
\end{figure}

\subsection{Use Case Diagram}

\begin{figure}[H]
  \centering
  \fbox{\parbox{0.85\textwidth}{\centering Placeholder: Use case diagram (select vulnerability → generate scenario → inspect node telemetry)}}
  \caption{Use case diagram for attack-path generation}
  \label{fig:use-case-diagram}
\end{figure}

\section{Design}
\subsection{Graph Data Model}

The attack graph is a simple, on-demand diagram that helps an analyst understand how an attacker could move through the system. It represents a generated scenario as a sequence of human-readable steps — each step has a short title, a clear description of the action, and a note about how serious it is. The diagram shows those steps as boxes connected by arrows to indicate possible flows or next actions. When an analyst clicks a box, the interface shows helpful details and suggested countermeasures for that step.

\subsection{Sequence Diagram (runtime)}

This sequence diagram summarizes the main runtime interaction for on-demand scenario generation: the Dashboard, the backend orchestration service, the AI gateway, and the frontend renderer. It visualizes the request flow (user selection → generation request → AI call → parsing/validation → interactive rendering) and highlights where failures (timeouts, malformed responses) are handled. The diagram emphasizes high-level handoffs rather than implementation details.

\begin{figure}[H]
  \centering
  \fbox{\parbox{0.85\textwidth}{\centering Placeholder: Sequence diagram (Dashboard → Controller → Service → AI → Parser → Dashboard)}}
  \caption{Sequence diagram for runtime scenario generation}
  \label{fig:attack-graph-sequence}
\end{figure}

\section{Implementation}

The implementation focuses on a compact orchestration layer that connects the Dashboard to an AI-based scenario generator while keeping the UI interactive and user-friendly.

- Backend: a lightweight service accepts generation requests, enriches them with vulnerability context, calls the configured AI gateway, and validates and normalizes the returned scenario. The service converts the scenario into an ephemeral graph optimized for client rendering and returns it to the Dashboard. It also handles error conditions (missing records, invalid AI output, gateway timeouts) and surfaces clear HTTP status codes and messages for the frontend to display.

- Frontend: the Dashboard triggers scenario generation for a selected vulnerability, displays progress and error states, and renders the returned graph for interactive inspection. Selecting a node reveals tactical details and suggested countermeasures in the intelligence pane; inventory filters allow analysts to narrow focus before generation.

- Operational notes: generated graphs are ephemeral for this sprint and are not persisted. Raw AI responses are retained only when necessary for debugging and must be protected in production. AI gateway settings (endpoint, model, timeout) are configurable via environment variables and should be tuned for reliability and operator monitoring.

\begin{figure}[H]
  \centering
  \fbox{\parbox{0.85\textwidth}{\centering Placeholder: UI screenshot of Dashboard attack graph rendering}}
  \caption{Dashboard attack graph (placeholder)}
  \label{fig:graph-output}
\end{figure}

\section{Conclusion}

This sprint establishes the attack-graph generation workflow: extraction from source artifacts, enrichment with vulnerability metadata, and rendering to visual formats. The produced graphs will be used in later chapters to illustrate attack paths and support evaluation.
